\documentclass[letterpaper,10pt]{article}
\usepackage[left=2cm,top=3cm,right=2cm,bottom=4cm]{geometry} 
\usepackage[utf8]{inputenc}
\usepackage[activeacute,spanish]{babel}
\usepackage{multirow}	% Para poder unir filas en las tablas
\usepackage{colortbl}	% Para colorear tablas
\usepackage{latexsym,amssymb,oldlfont,amsmath}
\usepackage{graphicx,graphics}
\usepackage{multicol}
\usepackage{color}
\usepackage{listings}
\usepackage{lastpage}
\usepackage{enumitem}



%\usepackage{fontspec}
%\setmainfont{Times New Roman}
\newcommand{\R}{{\mathbb R}}
\newcommand{\si}{{\mathcal S}^n}
\renewcommand{\theequation}{\thesection.\arabic{equation}} % Cambia formato salida numeracion equation
\makeatletter\@addtoreset {equation}{section} \makeatother % Tambien se usa
\renewcommand{\baselinestretch}{1.3}
\usepackage{amssymb}
\usepackage{hyperref}
\definecolor{midnightblue}{rgb}{0.1,0.1,0.44}


\newcommand{\dis}{\displaystyle}
\renewcommand{\theenumii}{\arabic{enumii}}

\def\Z{\mathbb{Z}}
\def\N{\mathbb{N}}
\def\R{\mathbb{R}}
\def\M{\mathcal{M}_n(\mathbb{R})}
\def\sin{\operatorname{sen}}
\def\e{\operatorname{e}}
\DeclareMathOperator{\arcsec}{arcsec}
\DeclareMathOperator{\arccot}{arccot}
\DeclareMathOperator{\arccsc}{arccsc}
\providecommand{\abs}[1]{\lvert#1\rvert}


\usepackage{etoolbox,fancyhdr,xcolor}
\fancyhf{} % sets both header and footer to nothing
\renewcommand{\headrulewidth}{0pt}
%\renewcommand{\footrulewidth}{0.5pt}
\pagestyle{fancy}
\headheight=60pt %para cambiar el tamaño del encabezadox
\fancyhead[L] %la "L" indica a la izquierda
{
\begin{minipage}{6cm}
\includegraphics[width=0.9\textwidth]{logo_udp.eps}
\end{minipage}
}

%\fancyhead[C]
%{
%\textbf{Cálculo II - Guía 1}
%}


%\fancyhead[R] %la "L" indica a la izquierda
%{

%}
%\begin{minipage}{4cm}
%\small %tamaño de letra pequeño
%{
%\textsf
%{
%Universidad de Chile\\
%...}}
%\end{minipage}}



%\fancyfoot[C]{}
%{
%\begin{minipage}{10cm}
%\color{midnightblue}
%\begin{center}Casa Matriz: Monjitas 527 Piso 15 - Santiago de Chile\\
%Fono (02)24367255 \end{center}
%\end{minipage}
%}
\fancyfoot[C]{}
%Página \thepage\ de \pageref{LastPage}}% Páginas Pares numeradas a la derecha
%\fancyfoot[LO]{\thepage} % Páginas Impares numeradas a la izquierda

\title{C2_Ay6}

\begin{document}

\begin{center}
Cálculo II\\
Ayudantía Nº6\\
Primer Semestre 2017
\end{center}

\begin{enumerate}
\item Determine si las siguientes afirmaciones son verdaderas o falsas y justifique su respuesta:

\begin{enumerate}[label=\emph{\alph*)}]

\item Si $h$ es una función continua y $f$, $g$ son derivables, entonces para la función:\\
\begin{center}
$\dis F(x) = \int_{f(x)}^{g(x)} h(t)\, dt$
\end{center} 
Se tiene que $F'(x)=h(g(x))g'(x)-h(f(x))g'(x)$

\item Si $f(x)$ es una función continua y positiva entonces:
\begin{center}
$\dis F(x) = \frac{\int_{0}^{x} t * f(t)\, dt}{\int_{0}^{x} f(t)\, dt}$
\end{center} 
Es creciente para $x \geq 0$

\item Si $f$ es continua y $\int_{0}^{4} f(x)\, dx = 40$, entonces $\int_{0}^{2} f(2x)\, dx = 5$
\end{enumerate}
\vspace{2mm}
\item Calcule:
$\displaystyle \lim_{x\to 0} \; {1\over x} \; \int ^x_0 (1 + \sin
2t)^{{1\over t}} dt$


\item Si $f$ es continua y $\displaystyle \int_0^9f(x)dx=4$, encuentre $\displaystyle \int_0^3xf(x^2)dx$

\item Calcule la siguiente integral definida: $\dis \int_{0}^{\pi/6} \frac{\sin(2x)}{\cos^2(2x)} \; dx$


\item Calcule la siguiente integral definida: $\dis \int_{1/2}^{\sqrt{3}/2} \frac{6}{\sqrt{1-t^2}} \; dt$





\end{enumerate}




\vspace{10mm}




\end{document}
